\chapter{Implicaciones}

\titular{El Paquete Económico 2027 mantiene el gasto plano en términos reales y recompone su
interior. \textbf{Lo que más cambia no es una cifra: es dónde vive el perímetro de exclusión.} La
cláusula de la Ley de Ingresos ---que se renegocia cada año--- desapareció, y una ley
permanente que reforma la LFPRH excluye ahora a las empresas públicas de otro cómputo, el
del gasto corriente estructural.}

\section{Cinco implicaciones, cada una con la cifra que la sostiene}

\subsection*{1. El techo no se mueve; la composición sí}

El gasto neto pagado queda en \textbf{$-$0,04~\% real} en la línea G y el gasto neto total en
\textbf{$+$0,3~\%} en la línea P. Dentro de ese techo, las \textbf{previsiones para gastos
obligatorios sin pensiones caen 4,0~\% real} y el \textbf{perímetro pensionario sube
3,9~\%}. \textbf{Un documento que compare totales concluirá que no pasó nada.}

La implicación para quien lea el presupuesto de 2027: \textbf{las preguntas interesantes están
todas por debajo del agregado}, y este año la fuente que permite bajar ---el analítico del
proyecto--- no estuvo disponible.

\subsection*{2. La regla fiscal cambió de domicilio legal}

El artículo 1o. de la iniciativa \textbf{no trae cláusula de exclusión ni tope}. La
reforma a la LFPRH publicada el 9 de abril de 2026 dispone que «el gasto de las empresas
públicas del Estado \textbf{no se contabilizará}» en la base del gasto corriente estructural,
y el agregado que la regla mide \textbf{cae 51,3~\% en pesos corrientes} mientras su límite
pasa de \textbf{10,9 a 5,2~\% del PIB}.

\textbf{No es una implicación sobre disciplina fiscal, es una sobre gobernanza:} un perímetro
que vivía en una ley anual, sujeta a la Cámara cada septiembre, ahora vive en una ley
permanente. \textbf{Y el renglón que permitiría medir la regla ---el balance sin
inversión--- lleva dos ejercicios sin publicarse.}

\subsection*{3. Un solo programa explica el movimiento más grande}

El \textbf{Ramo 18 Energía pierde 181.804,9 mdp} y el programa «Articulación de la Política
de Hidrocarburos» pierde \textbf{182.376,0}. La diferencia, \textbf{$+$571,1 mdp}, es lo que
gana el resto del ramo.

La implicación es de método: \textbf{la caída de un ramo casi nunca es la caída de la
política que su nombre sugiere}, y en este caso la fuente misma lo dijo. Un lector que
hubiera titulado «recorte de 69~\% a la energía» habría descrito una operación financiera con
Pemex.

\subsection*{4. Por persona, las pensiones no crecen}

El perímetro pensionario sube \textbf{3,9~\% real} y la población de 65 años y más
\textbf{4,29~\%}: \textbf{por persona de esa edad el gasto cae 0,4~\% real}. En educación
ocurre lo contrario: el FONE sube 6,8~\% real y la población en edad escolar básica cae
0,95~\%, de modo que \textbf{por persona sube 7,8~\%}.

La implicación: \textbf{los montos y las razones por persona dan lecturas opuestas}, y las
dos son correctas. Cuál se use depende de la pregunta, y \textbf{la pregunta hay que
declararla.}

\subsection*{5. La fuente oficial deflacta con el índice equivocado}

El CGPE declara un deflactor del PIB de \textbf{4,04~\%} y sus cuadros presupuestales están
calculados con \textbf{1,0325}, la inflación promedio del INPC. Despejado de nueve renglones
independientes.

La implicación es doble. Para el lector: \textbf{las variaciones reales de este documento y
las del CGPE no son comparables sin ajustar}, y la diferencia vale unos 0,8 puntos. Para el
género: \textbf{el error que este proyecto persigue desde 2020 ---mezclar convenciones de
deflactación--- lo comete la fuente primaria}, en el mismo documento que declara el
deflactor correcto tres páginas antes.

\section{Lo que este documento no pudo hacer, y a quién le toca}

\begin{list}{}{\setlength{\leftmargin}{1.4em}\setlength{\itemsep}{2pt}}
\item[$\bullet$] \textbf{El diff institucional y todo el análisis por programa}, por el 404 de
los cuatro analíticos del proyecto. \textbf{No es una limitación de método ni de tiempo: es
una fuente que el servidor de la Secretaría no entregó el día de la entrega.}
\item[$\bullet$] \textbf{La auditoría de etiquetado de los anexos transversales}, por el 503
de la base de datos abiertos. Es la aportación más específica del método de este proyecto y
se quedó sin hacer.
\item[$\bullet$] \textbf{La descomposición del cambio del acervo de deuda}, que no pasó su
criterio de aceptación y por eso no se publica.
\item[$\bullet$] \textbf{La exposición de motivos del PPEF}, que el portal enumera y no
sirve, por sexto ejercicio consecutivo.
\end{list}

\section{Una implicación sobre las fuentes, que no es menor}

De las tres paradas que la instrucción fija para obtener los analíticos ---portal de la
Secretaría, portal de datos abiertos y Transparencia Presupuestaria--- \textbf{las tres
fallaron el día de la entrega}, cada una de manera distinta: el árbol del ejercicio 2027 no
existía, las ranuras del portal de datos abiertos seguían apuntando al ejercicio anterior, y
Transparencia Presupuestaria devolvió \textbf{1.923 bytes de cáscara sin un solo enlace a
archivo}, la misma cifra que en 2024.

\textbf{Un paquete económico que se publica sin sus analíticos es un paquete que no se puede
auditar el día que se presenta}, que es justamente el día en que la discusión pública lo
necesita. Este documento se escribió a nivel de ramo y lo declaró en cada capítulo; no todo
lector de un paquete económico tiene esa disciplina, y \textbf{la ausencia de la fuente
empuja al género entero hacia la comparación de totales}, que es donde no pasa nada.

\begin{ficha}{implicaciones}
\fila{Perímetro}{Ninguno propio: este capítulo no introduce cifras nuevas. Cada cifra que
cita remite al capítulo donde se estableció, con su perímetro declarado ahí.}
\fila{Línea de comparación}{La que corresponda al capítulo de origen, etiquetada en cada
mención.}
\fila{Deflactor}{Del PIB, 1,040, en todas las variaciones reales citadas.}
\fila{Año de los pesos}{Pesos de 2027.}
\fila{PIB y añada}{PIB de 2027 del CGPE 2027; el de 2026, el aprobado.}
\fila{Denominadores}{Los del capítulo de horizonte demográfico, declarados ahí.}
\fila{Fuentes}{Las de los capítulos citados. No hay fuente exclusiva de este capítulo.}
\fila{Qué no puede afirmarse con ellas}{Ninguna de estas implicaciones es una recomendación
de política. Describen lo que el paquete hace y lo que su documentación permite o impide
verificar; qué debería hacerse con ello es una decisión que este documento no toma.}
\end{ficha}
